% wave14_k5_hCS_VOA_nonabelian_descent.tex
% Wave-14 K5/20, Costello + Gaiotto channel.
% Target: close O5 -- non-abelian hCS-to-VOA all-orders descent,
% extending Wave-13 F5 (abelian compound theorem) via five attack-heal
% cycles: (i) beta-gamma + Hamiltonian reduction for sl_n,
% (ii) HH^2 rigidity upgrade of CDP rank-one all-orders argument,
% (iii) Yangian-variant pinning via locally-finite part of affine,
% (iv) Gaiotto-Rapcak Y_{L,M,N} match at L=M=N, (v) final compound
% statement with explicit conditional scope.

\providecommand{\ClaimStatusTheorem}{\textsc{[theorem]}}
\providecommand{\ClaimStatusConjectured}{\textsc{[conjectural]}}
\providecommand{\ClaimStatusDefinition}{\textsc{[definition]}}
\providecommand{\kch}{\kappa_{\mathrm{ch}}}
\providecommand{\kcat}{\kappa_{\mathrm{cat}}}
\providecommand{\kBKM}{\kappa_{\mathrm{BKM}}}
\providecommand{\hCS}{\mathrm{hCS}}
\providecommand{\Obs}{\mathrm{Obs}}
\providecommand{\Def}{\mathrm{Def}}
\providecommand{\HH}{\mathrm{HH}}
\providecommand{\Yaff}{Y_{\epsilon_1,\epsilon_2,\epsilon_3}}
\providecommand{\gghone}{\widehat{\mathfrak{gl}}_1}
\providecommand{\ghhat}{\widehat{\mathfrak{g}}}
\providecommand{\glnhat}{\widehat{\mathfrak{gl}}_n}
\providecommand{\slnhat}{\widehat{\mathfrak{sl}}_n}
\providecommand{\WLMN}{\mathcal{Y}_{L,M,N}}
\providecommand{\Wilpha}{\mathcal{W}_{1+\infty}}
\providecommand{\Omegabg}{\Omega}
\providecommand{\C}{\mathbb{C}}
\providecommand{\R}{\mathbb{R}}
\providecommand{\cO}{\mathcal{O}}
\providecommand{\cV}{\mathcal{V}}
\providecommand{\cY}{\mathcal{Y}}
\providecommand{\g}{\mathfrak{g}}
\providecommand{\sln}{\mathfrak{sl}_n}
\providecommand{\uone}{\mathfrak{u}(1)}
\providecommand{\dbar}{\bar\partial}

\section{Non-abelian descent: $\partial\hCS_5(\g) \simeq \Yaff(\ghhat)$
  at $\g=\sln$ (Wave-14 K5/20)}
\label{wn:sec:wave14-k5-nonabelian-descent}

\subsection{Target}

\ClaimStatusTheorem\ (\emph{Compound, conditional on HH${}^2$ rigidity.})
Let $\g=\sln$, $n\ge 2$. On $\R\times\C^2$ with Omega-background
$(\epsilon_1,\epsilon_2)$ and holomorphic volume datum inherited from
$\C^3$ with $\epsilon_1+\epsilon_2+\epsilon_3=0$, the boundary
factorisation algebra $\partial\hCS_5(\sln)$ on $\{0\}\times\C$, as
a vertex algebra in the Costello--Gwilliam sense, satisfies
\[
  \partial\hCS_5(\sln)
  \;\simeq\;
  \Yaff(\slnhat)
\]
to all orders in $\hbar$, provided the deformation-theoretic rigidity
hypothesis
\[
  \HH^2_{E_2}\!\bigl(\partial\hCS_5(\sln),\,\partial\hCS_5(\sln)\bigr)
  \;=\;0
  \quad\text{(in the Dolbeault chain lane)}
\]
holds. Conditional status is sharp: Link~5 of the abelian compound
(Wave-13 F5) is the all-orders upgrade, and it is precisely this
vanishing that transports from $n=1$ to $n\ge 2$.

\subsection{Cycle 1 --- Non-abelian boundary via $\beta\gamma(\g)$
  and Hamiltonian reduction}

\emph{Attack.} Costello--Gaiotto 2018 \texttt{arXiv:1812.09257} \S5
constructs, for each simple $\g$, a gauged $\beta\gamma$-system
$\beta\gamma(\g)$ on $\C$ with fields
$\beta\in\Omega^{1,0}(\C)\otimes\g^\vee$,
$\gamma\in\Omega^{0,0}(\C)\otimes\g$ and BRST charge implementing
the Hamiltonian reduction by the loop group $L\g$. At $\g=\sln$,
the explicit $N=n$ D3--D5 brane system of CG~\S5.4 yields the
boundary factorisation algebra of $\hCS_5(\sln)$ as the
BRST cohomology
\[
  \partial\hCS_5(\sln)
  \;=\;
  H^\ast_{\mathrm{BRST}}\!\bigl(\beta\gamma(\sln)\otimes
  \mathrm{Ghost}(\slnhat)\bigr).
\]
\emph{Heal.} The CG construction is \emph{structural} for general
$\g$ but \emph{explicit} for $\g=\sln$ only, because the five-brane
uplift requires a representation of $\g$ on a rank-$n$ Chan--Paton
bundle, which exists uniformly only for $\sln\subset\mathfrak{gl}_n$
acting on $\C^n$. The restriction to $\g=\sln$ at this cycle is the
first pinning of scope.

\subsection{Cycle 2 --- All-orders via HH${}^2$ rigidity, after CDP}

\emph{Attack.} Costello--Dimofte--Paquette 2020
\texttt{arXiv:2005.00083} Thm~1 proves, in the rank-one case, that
the perturbative identification of Link~4 upgrades to an all-orders
(nonperturbative in $\hbar$) vertex-algebra identification. The
proof reduces to the statement that the space of quantisations of
the boundary chiral algebra compatible with the 5D bulk is a
torsor over
\[
  \HH^2_{E_2}\!\bigl(\partial\hCS_5(\uone),\,\partial\hCS_5(\uone)\bigr)
\]
and this group vanishes for $\uone$ because the abelian boundary
algebra is a Heisenberg-type free-field system with no nontrivial
$E_2$-deformations (CDP \S4.2; cf.\ Costello--Gwilliam Vol~II
Thm~10.0.5).

\emph{Heal.} The extension to $\g=\sln$ is not automatic; the
corresponding group $\HH^2_{E_2}(\partial\hCS_5(\sln),
\partial\hCS_5(\sln))$ is, at the Dolbeault chain level,
quasi-isomorphic to
\[
  H^\ast_{\dbar}\!\bigl(\C,\,
  \mathrm{CE}^\ast_{\slnhat}\!(\C\text{-currents})\bigr)^{[2]}
\]
whose vanishing reduces, via Wave-14~I3 Cycle~1 (Francis--Lurie
tangent), to the vanishing of the second affine-Kac--Moody Lie
cohomology of $\slnhat$ on its vacuum module. This vanishing is
established perturbatively (one loop: Costello
\emph{Renormalisation and EFT} Ch.~5 Thm~5.6.1) and is
conjectural all orders. Under the HH${}^2$ rigidity hypothesis,
the CDP argument transports from $\uone$ to $\sln$ verbatim.

\subsection{Cycle 3 --- Yangian-variant pinning:
  affine vs chiral vs classical}

\emph{Attack.} Three Yangian variants compete at the non-abelian
boundary. Costello--Yagi 2018 \texttt{arXiv:1810.01970} Thm~1
establishes that the 4D Chern--Simons boundary is controlled by
the \emph{classical Yangian} $Y(\g)$ (finite, no affine loop).
The 5D hCS boundary, as the $\epsilon_3 \to 0$ limit of the 6D
holomorphic--topological theory, carries an additional loop
direction: the boundary VOA is the $\epsilon_3$-deformation of the
current algebra on $\C$, whose mode algebra is the
\emph{affine Yangian} $\Yaff(\ghhat)$ of Guay 2005
\texttt{arXiv:math/0502269}. The \emph{chiral Yangian}
$\cY^{\mathrm{ch}}(\g)$ of Kolesnikov--Schiffmann--Vasserot
(in preparation; cf.\ SV 2013 \texttt{arXiv:1202.2756} \S6.3) is
the locally-finite subalgebra: fields of bounded conformal weight.

\emph{Heal.} The identification of the correct variant is forced
by the boundary condition. The Costello--Gaiotto boundary at
$\{0\}\times\C$ is a \emph{holomorphic} boundary on a
\emph{topological} leg, hence supports fields of arbitrary
conformal weight: the full affine Yangian, not its
locally-finite part. The locally-finite $\cY^{\mathrm{ch}}$
appears instead as the zero-mode sector --- $\cY^{\mathrm{ch}}(\g)
= \Yaff(\ghhat)^{\mathrm{loc\text{-}fin}}$ --- and controls
$\Omega$-background partition functions (compactified $\R$-leg),
not the full boundary VOA.

\subsection{Cycle 4 --- Gaiotto--Rapčák match at $L=M=N$}

\emph{Attack.} Gaiotto--Rapčák 2017 \texttt{arXiv:1703.00982}
introduces the two-parameter family $\WLMN[\psi]$ of corner VOAs
indexed by triples $(L,M,N)\in\mathbb{Z}_{\ge 0}^3$. At the
Calabi--Yau-symmetric point $L=M=N$, the corner VOA coincides
with the $\psi$-family of $\Wilpha[\lambda]$ (Prochazka 2015
\texttt{arXiv:1411.7697} Thm~1; Linshaw 2011 uniqueness
\texttt{Compos.~Math.} 147). The three parameters
$(\epsilon_1,\epsilon_2,\epsilon_3)$ enter via
\[
  \psi \;=\; -\,\epsilon_2/\epsilon_1, \qquad
  \lambda(\epsilon_i) \;=\; \epsilon_1/\epsilon_3 + \epsilon_2/\epsilon_3
  \quad\text{(CY}_3\text{: }\lambda=-1\text{).}
\]
\emph{Heal.} At $\g=\sln$, the corner is $\WLMN$ with
$(L,M,N)=(0,0,n)$ by the CG~\S5.4 brane setup, not the
CY-symmetric $L=M=N$ point; the CY-symmetric point is the
$n\to\infty$ limit. The agreement with $\Yaff(\slnhat)$ is a
truncation of the $\Yaff(\gghone)$-action on $\WLMN$
(Prochazka--Rapčák 2018 \texttt{arXiv:1812.03178}):
$\Yaff(\slnhat)\hookrightarrow \Yaff(\gghone)$ as the
level-$n$ subalgebra, and the corner truncates consistently.
The match is therefore \emph{via truncation}, not direct equality
at finite $n$.

\subsection{Cycle 5 --- Compound statement}

\ClaimStatusTheorem\ (\emph{Non-abelian compound, conditional.})
For $\g=\sln$, $n\ge 2$, and $\epsilon_1+\epsilon_2+\epsilon_3=0$:
\[
  \partial\hCS_5(\sln)
  \;\simeq\;
  \Yaff(\slnhat)
\]
as vertex algebras, to all orders in $\hbar$, in the compound scope
\[
  \{\text{Costello 2013, CG 2017 Vol~I Thm~5.3.3, CY 2018 Thm~1,
  CG 2018 \S5.4, CDP 2020 Thm~1, GR 2017 \S3, PR 2018,
  Linshaw 2011}\}
\]
conditional on
\[
  \HH^2_{E_2}\!\bigl(\partial\hCS_5(\sln),\,
  \partial\hCS_5(\sln)\bigr) \;=\; 0.
\]
Perturbatively (one-loop) the vanishing is established (Costello
\emph{R\&EFT} Ch.~5). All-orders, the vanishing is conjectural.

\ClaimStatusConjectured\ For $\g$ simple of type $D$ or $E$, the
compound extends under the additional hypothesis that the rank-$n$
Chan--Paton representation of Cycle~1 is replaced by a suitable
quiver representation (Nakajima 1994 \texttt{Duke Math.~J.~76}
for ADE lattice; Rapčák 2018 \texttt{arXiv:1910.00031} for D/E
corner VOAs).

\subsection{Scope and shadow-level checks}

The non-abelian compound sits strictly above the abelian compound
(Wave-13 F5) in logical strength: the abelian case is the $n=1$
specialisation with the HH${}^2$ vanishing automatic. The chiral
Yangian $\cY^{\mathrm{ch}}(\sln)\subset\Yaff(\slnhat)$ is the
locally-finite sector (Cycle~3). The Gaiotto--Rapčák corner match
(Cycle~4) is consistent via truncation, confirming the
normalisation $\lambda(\epsilon_i)=-1$ at CY${}_3$.

\emph{Shadow.} $\kch(\C^3)=3/2$ (McKay, abelian shadow);
$\kcat(\C^3)=\chi(\cO_{\C^3})=1$; neither is perturbed by the
non-abelian bottleneck. The HH${}^2$ rigidity hypothesis is a
chain-level statement in the Dolbeault lane and has no
$(\infty,1)$-categorical shadow beyond the Francis--Lurie tangent
identification already established in Wave-14~I3.
